% Outline: Related Work (1 page cap). Ryoo craft: pure gap-framing, never survey;
% batch-cites to establish each field; bold run-in headers; differentiate by AXIS;
% first-person "Our work ..." differentiators; concede-then-pivot, never dismiss.
\section{Related Work}

\subsection{Learned Length Prediction for Scheduling}
A substantial line learns output-length signals to approximate
shortest-job-first serving~\cite{ssjf2024,fuLTR2024,tao2026ranking,elis2025},
progressing from length-bucket classification to relative-rank training, the
insight we inherit: scheduling needs \emph{order}, not absolute
length~\cite{fuLTR2024}. PARS, now peer-reviewed at ISC High Performance
2026, pairs a BERT encoder with pairwise margin training and demonstrates
cross-model robustness~\cite{tao2026ranking}. However, every predictor in
this line reads the prompt or scalar request statistics; none consumes
tool-schema text, and none is evaluated on requests whose tools were absent
from training. Our work fills this gap: schema text as
predictor input, and a novelty-stratified evaluation in which 95.5\% of test
requests carry tool combinations unseen in training --- the regime our
single-client live capture exhibits throughout.

\subsection{Prediction-Free Scheduling}
A newer position holds that tail latency does not need prediction at all.
UniBoost's distribution-aware boosting beats even an \emph{oracle}-length
SRPT baseline by 35.1\% on P99 time-to-last-token in its Table~2
configuration~\cite{beyondPrediction2026}, and starvation-free preemptive
designs make adjacent arguments~\cite{fastServe2026,tie2026}. We
therefore make no tail claims: if an oracle loses on P99, better ranking
accuracy will not repair that, so our endpoints are means and goodput. The concession is scoped by the papers' own designs, however:
UniBoost evaluates no tool-calling workload, disables prefix caching for
experimental control, and is prediction-free but not feature-free --- its
scheduler consumes prompt length, the same class of signal our weakest
baseline uses. Whether its result transfers to agentic traffic is, by its
own construction, open.

\subsection{Scheduling Under Imprecise Information}
JITServe is the strongest published treatment of acting on unreliable
estimates: quantile-forest \emph{upper bounds} on length, refined every
$\sim$50 generated tokens, drive grouped goodput scheduling within 3--9\% of
an oracle on a 16-GPU cluster --- with BERT-class predictors explicitly
rejected at 56--187\,ms under load against a forest costing 11--24\,ms at
those same loads (its oft-quoted 7\,ms is the lightest, 8-requests/s
point)~\cite{jitServe2026}. Its tool awareness is structural: tools enter as
identity and duration in a dependency graph, never as text into the length
model. We position our gate as complementary, not competitive. JITServe
absorbs uncertainty by always scheduling against a pessimistic bound; our
gate \emph{detects} where trust has no measured basis and abstains,
preserving arrival order --- the ``explicit fallback policies'' its authors
name as an open problem under distribution shift. Agent-serving systems that
exploit workflow structure~\cite{agentix2026,infercept2024} are orthogonal:
they schedule around tool-call \emph{interruptions}, not around length
uncertainty.

\subsection{Schemas as Model Input; Unseen-Tool Evaluation}
Fine-tuning small encoders on schema text is established practice in tool
learning --- it is how function-calling models are
trained~\cite{toolace2024,qin2024toolllm,patil2024gorilla} --- and
Switchcraft packs function signatures into a DistilBERT to route requests
across models by predicted \emph{correctness}~\cite{switchcraft2026}. The
distinction we claim is therefore narrow and specific: schema text
$\rightarrow$ output-length rank $\rightarrow$ queue order is, to our
knowledge, unoccupied; Switchcraft itself lists benchmark-shift
quantification as future work; our Cold-Start Transfer strata measure that
axis. For the stratified design we follow tool-learning precedent:
ToolLLM evaluates at unseen-instruction, unseen-tool, and unseen-category
levels~\cite{qin2024toolllm}, licensing our S1--S4 \emph{design} though not
our metric (it scores task success, not length rank). BFCL's crowd-sourced
split targets contamination control, a different axis~\cite{patil2025bfcl},
and Gorilla's holdout is random --- its ``zero-shot'' denotes retriever
absence, not unseen APIs~\cite{patil2024gorilla} --- so neither substitutes
for tool-novelty stratification.

\subsection{The Serving Substrate: Gateways, Admission, Prefix Reuse}
Our contribution sits between two systems layers that are usually studied
separately. Below the engine, continuous
batching~\cite{orca2022} and paged attention~\cite{vllm2023} set the
execution substrate, and prefix reuse --- RadixAttention in
SGLang~\cite{sglang2024} and the prefix caching now standard in
vLLM --- exploits the repetition our capture exhibits, a
schema resent byte-identically on every turn of a given call kind; our
ordering runs hold it
disabled to isolate ordering, and one round enables it to measure it. Above the engine, gateways
own admission: overload control in production RPC systems sheds load
by priority when queues exceed bounds~\cite{dagor2018}, and the
gateway we deploy on implements the same pattern as queue guards with
priority bypass. Admission decides \emph{whether and when} a request
reaches the engine; the engine's scheduler decides \emph{which}
waiting request runs next. Our predictor feeds the second decision,
and its budget is set by the first, so a scoring path
slower than the configured admission timeout is excluded regardless of
its ranking accuracy; the timeout itself is an operator choice.

